\section{Background and Related Work}
\textbf{Long-term memory (LTM).} 
Persistent LTM is crucial for LLM-based agents operating over extended horizons~\citep{wang2025karma, li2025hello}.
Recent work has explored diverse architectural designs for modeling LTM. 
LangMem~\citep{langmem2025} provides a modular framework that supports multiple memory types, while A-Mem~\citep{xu2025mem} adopts a Zettelkasten-inspired design that links structured knowledge units to facilitate consolidation. 
Mem0~\citep{chhikara2025mem0} proposes a scalable extract-update pipeline and extends it to a graph-based variant for structured reasoning, and Zep~\citep{rasmussen2025zep} represents memory as a temporal knowledge graph to enable cross-session and time-aware reasoning. Although effective in organizing and retrieving information, these approaches largely rely on predefined memory structures or heuristic update rules. 
As memory grows, such designs commonly suffer from increased system complexity and lack adaptive, learning-based strategies for prioritization and forgetting.
In contrast, our work aims to learn an adaptive memory policy that allows agents to dynamically decide what to store, update, or forget, depending on task demands and long-term utility.
\\
\textbf{Short-term memory (STM).} 
STM in agentic LLMs primarily concerns context selection and retrieval~\citep{wang2024agent, jin2024llm}.
Retrieval-Augmented Generation (RAG)~\citep{pan2025memory, salama2025meminsight, kagaya2024rap} is the dominant paradigm, expanding usable context by injecting retrieved content into prompts.
While effective, RAG does not fundamentally prevent context explosion in long-horizon settings and may introduce irrelevant or distracting information.
To address this issue, ReSum~\citep{wu2025resum} periodically compresses interaction histories into compact reasoning states, allowing agents to operate beyond fixed context-window constraints. 
Yet its summarization schedule remains largely predefined, and aggressive compression risks discarding rare but crucial details.
Our approach instead enables agents to \emph{learn} when and how to retrieve, summarize, or filter context, achieving a more flexible balance between efficiency and information preservation.
\\
\textbf{Reinforcement learning for LLMs.} 
Reinforcement learning has become an effective paradigm for improving the decision-making and reasoning capabilities of LLM-based agents~\citep{yao2022react, jin2025search, qian2025toolrl, chaudhari2025rlhf}. 
Among recent advances, GRPO~\citep{shao2024deepseekmath} enhances stability by optimizing policies based on the relative quality of sampled trajectories, removing the need for an explicit value function. 
GRPO and its variants~\citep{gilabert2025terminology, wang2025grpo} have shown strong performance in complex reasoning tasks.
However, existing RL-based systems generally treat memory as a static or external component, making them ill-suited for the discontinuous and fragmented trajectories associated with memory operations~\citep{yan2025memory, zhang2025memory}. 
In contrast, our work integrates RL directly into the memory management process, enabling unified training of both language generation and memory operations.


